\section{Рабочий проект}

Программно-информационная система обнаружения дефектов дорожного покрытия включает несколько программных модулей: модуль запуска приложения, главный модуль приложения, модуль обнаружения дефектов, модуль работы с файлами, модуль галереи изображений, модуль просмотра изображений, модуль управления визуальным оформлением.


\subsection{Модуль "<main.py">}

Модуль "<main.py"> является точкой входа в программно-информационную систему обнаружения дефектов дорожного покрытия. Основной задачей модуля является создание экземпляра приложения Qt, инициализация главного окна программы и запуск цикла обработки событий пользовательского интерфейса.

Функции модуля запуска приложения представлены в таблице \ref{main:table}.

\renewcommand{\arraystretch}{0.8}
\begin{xltabular}{\textwidth}{|p{3.5cm}|p{3.5cm}|>{\setlength{\baselineskip}{0.7\baselineskip}}X|}
	\caption{Описание функции, используемой в модуле запуска приложения\label{main:table}}\\
	\hline \centrow \setlength{\baselineskip}{0.7\baselineskip} Название функции & \centrow \setlength{\baselineskip}{0.7\baselineskip} Аргументы & \centrow Описание функции \\ 
	\hline \centrow 1 & \centrow 2 & \centrow 3\\ \hline
	\endfirsthead
	\caption*{Продолжение таблицы \ref{main:table}}\\
	\hline \centrow 1 & \centrow 2 & \centrow 3\\ \hline
	\finishhead
	main & Аргументов нет & Точка входа в приложение. Выполняет создание объекта класса QApplication, создание экземпляра класса MainWindow, отображение главного окна программы и запуск основного цикла обработки событий пользовательского интерфейса. \\ \hline	
\end{xltabular}
\renewcommand{\arraystretch}{1.0}

\subsection{Модуль "<main\_window.py">}

Модуль "<main\_window.py"> содержит класс "<MainWindow">, реализующий главное окно программно-информационной системы обнаружения дефектов дорожного покрытия. Класс обеспечивает взаимодействие пользователя с системой, управление обработкой изображений, отображение результатов обнаружения дефектов, настройку параметров работы приложения и сохранение результатов обработки.

\subsubsection{Класс "<MainWindow">}

Методы класса "<MainWindow"> представлены в таблице \ref{mainwindow:table}.

\renewcommand{\arraystretch}{0.8}
\begin{xltabular}{\textwidth}{|p{5.5cm}|p{3.5cm}|>{\setlength{\baselineskip}{0.7\baselineskip}}X|}
	\caption{Описание методов класса "<MainWindow">\label{mainwindow:table}}\\
	\hline \centrow \setlength{\baselineskip}{0.7\baselineskip} Название метода & \centrow \setlength{\baselineskip}{0.7\baselineskip} Аргументы & \centrow Описание метода \\
	\hline \centrow 1 & \centrow 2 & \centrow 3\\ \hline
	\endfirsthead
	\caption*{Продолжение таблицы \ref{mainwindow:table}}\\
	\hline \centrow 1 & \centrow 2 & \centrow 3\\ \hline
	\finishhead
	
	\_\_init\_\_ & Аргументов нет & Конструктор класса. Выполняет инициализацию главного окна приложения, создание необходимых объектов и настройку пользовательского интерфейса. \\ \hline
	
	init\_window & Аргументов нет & Настройка основных параметров главного окна приложения, включая заголовок и размеры окна. \\ \hline
	
	init\_settings & Аргументов нет & Загрузка пользовательских настроек приложения и настройка каталога для хранения временных результатов обработки. \\ \hline
	
	init\_state & Аргументов нет & Инициализация внутренних структур данных и переменных состояния приложения. \\ \hline
	
	init\_managers & Аргументов нет & Создание экземпляров классов ThemeManager, FileManager и GalleryManager. \\ \hline
	
	init\_detector & Аргументов нет & Создание экземпляра класса Detector и загрузка обученной модели обнаружения дефектов дорожного покрытия. \\ \hline
	
	init\_menu & Аргументов нет & Формирование главного меню приложения и привязка действий пользователя к соответствующим обработчикам. \\ \hline
	
	init\_wid & Аргументов нет & Создание и настройка элементов пользовательского интерфейса главного окна приложения. \\ \hline
	
	init\_statusbar & Аргументов нет & Создание строки состояния и элементов отображения статистики обработки изображений. \\ \hline
	
	update\_progress & current, total, text & Обновление индикатора выполнения операции и отображение текущего состояния процесса. \\ \hline
	
	show\_welcome\_screen & Аргументов нет & Отображение стартового экрана приложения с краткой инструкцией по работе с системой. \\ \hline
	
	clear\_welcome\_screen & Аргументов нет & Удаление стартового экрана из рабочей области приложения. \\ \hline
	
	update\_results\_dir\_display & Аргументов нет & Обновление отображаемого пути к каталогу временных результатов обработки. \\ \hline
	
	set\_images & files & Сохранение списка выбранных изображений и запуск формирования миниатюр входных данных. \\ \hline
	
	save\_results & Аргументов нет & Сохранение обработанных изображений и отчёта в выбранный пользователем каталог. \\ \hline
	
	open\_processed\_results & Аргументов нет & Загрузка ранее сохранённых результатов обработки изображений. \\ \hline
	
	toggle\_input\_images & checked & Управление отображением панели входных изображений. \\ \hline
	
	set\_conf\_threshold & value & Установка порога уверенности обнаружения объектов нейронной сетью. \\ \hline
	
	sort\_results & mode & Сортировка результатов обработки по имени файла или количеству обнаруженных объектов. \\ \hline
	
	resizeEvent & event & Обработка изменения размеров окна приложения и перестроение области отображения результатов. \\ \hline
	
	select\_image & index & Выбор изображения для последующего просмотра в отдельном окне. \\ \hline
	
	open\_viewer & index & Открытие окна просмотра исходных изображений. \\ \hline
	
	open\_processed\_viewer & index & Открытие окна просмотра обработанных изображений. \\ \hline
	
	detect\_objects & Аргументов нет & Выполнение обнаружения дефектов дорожного покрытия на выбранных изображениях с использованием нейросетевой модели. \\ \hline
	
	closeEvent & event & Обработка завершения работы приложения и контроль сохранения результатов обработки. \\ \hline	
\end{xltabular}
\renewcommand{\arraystretch}{1.0}


\subsection{Модуль "<detector.py">}

Модуль "<detector.py"> содержит класс "<Detector">, реализующий обнаружение дефектов дорожного покрытия на изображениях с использованием предварительно обученной нейросетевой модели. Модуль обеспечивает загрузку изображений, подготовку данных для обработки, запуск нейросетевого анализа, визуализацию результатов обнаружения и сохранение обработанных изображений.

\subsubsection{Класс "<Detector">}

Методы класса "<Detector"> представлены в таблице \ref{detector:table}.

\renewcommand{\arraystretch}{0.8}
\begin{xltabular}{\textwidth}{|p{3.5cm}|p{3.5cm}|>{\setlength{\baselineskip}{0.7\baselineskip}}X|}
	\caption{Описание методов класса "<Detector">\label{detector:table}}\\
	\hline \centrow \setlength{\baselineskip}{0.7\baselineskip} Название метода & \centrow \setlength{\baselineskip}{0.7\baselineskip} Аргументы & \centrow Описание метода \\
	\hline \centrow 1 & \centrow 2 & \centrow 3\\ \hline
	\endfirsthead
	\caption*{Продолжение таблицы \ref{detector:table}}\\
	\hline \centrow 1 & \centrow 2 & \centrow 3\\ \hline
	\finishhead
	
	\_\_init\_\_ & model, conf\_threshold, max\_dim &
	Конструктор класса. Выполняет сохранение экземпляра нейросетевой модели, установку порога обнаружения объектов и максимального размера обрабатываемого изображения. \\ \hline
	
	set\_conf & value &
	Установка нового значения порога уверенности обнаружения объектов. \\ \hline
	
	process\_images & images, results\_dir, progress\_callback, stop\_flag &
	Выполнение обнаружения дефектов дорожного покрытия на наборе изображений. Выполняет загрузку и масштабирование изображений, запуск нейросетевой модели, подсчёт обнаруженных объектов, визуализацию результатов обнаружения и сохранение обработанных изображений. Возвращает список обработанных файлов, количество обнаруженных объектов на каждом изображении и общее число обнаруженных объектов. \\ \hline
	
\end{xltabular}
\renewcommand{\arraystretch}{1.0}


\subsection{Модуль "<file\_manager.py">}

Модуль "<file\_manager.py"> содержит класс "<FileManager">, предназначенный для выполнения операций с файлами и каталогами. Модуль обеспечивает загрузку исходных изображений, сохранение результатов обработки, формирование текстового отчёта, управление временными файлами, а также загрузку ранее сохранённых результатов обработки.

\subsubsection{Класс "<FileManager">}

Методы класса "<FileManager"> представлены в таблице \ref{filemanager:table}.

\renewcommand{\arraystretch}{0.8}
\begin{xltabular}{\textwidth}{|p{4.8cm}|p{3.5cm}|>{\setlength{\baselineskip}{0.7\baselineskip}}X|}
	\caption{Описание методов класса "<FileManager">\label{filemanager:table}}\\
	\hline \centrow \setlength{\baselineskip}{0.7\baselineskip} Название метода & \centrow \setlength{\baselineskip}{0.7\baselineskip} Аргументы & \centrow Описание метода \\
	\hline \centrow 1 & \centrow 2 & \centrow 3\\ \hline
	\endfirsthead
	
	\caption*{Продолжение таблицы \ref{filemanager:table}}\\
	\hline \centrow 1 & \centrow 2 & \centrow 3\\ \hline
	\finishhead
	
	\_\_init\_\_ & window &
	Конструктор класса. Сохраняет ссылку на главное окно приложения для последующего взаимодействия с его элементами интерфейса. \\ \hline
	
	open\_images & Аргументов нет &
	Открытие диалогового окна выбора изображений и передача выбранных файлов в главное окно приложения. \\ \hline
	
	save\_results & Аргументов нет &
	Сохранение обработанных изображений в выбранный каталог. Выполняет проверку существования файлов, обработку конфликтов имён и формирование отчёта по результатам обработки. \\ \hline
	
	create\_report & folder &
	Создание текстового отчёта о выполненной обработке изображений в каталоге \texttt{folder}. В отчёт включаются статистические данные и информация по каждому обработанному изображению. \\ \hline
	
	select\_results\_folder & Аргументов нет &
	Выбор каталога для хранения временных результатов обработки и сохранение выбранного пути в настройках приложения. \\ \hline
	
	clear\_results\_folder & Аргументов нет &
	Очистка каталога временных результатов обработки после подтверждения пользователем. \\ \hline
	
	open\_processed\_results & Аргументов нет &
	Открытие ранее сохранённых результатов обработки. Выполняет выбор каталога и загрузку информации из отчёта. \\ \hline
	
	\_load\_report & folder, report\_path &
	Чтение и анализ файла отчёта, расположенного по пути \texttt{report\_path}. Формирует структуру данных, содержащую сведения о результатах обработки изображений. \\ \hline
	
\end{xltabular}
\renewcommand{\arraystretch}{1.0}


\subsection{Модуль "<gallery\_manager.py">}

Модуль "<gallery\_manager.py"> содержит класс "<GalleryManager">, предназначенный для отображения изображений в пользовательском интерфейсе программы. Модуль обеспечивает формирование галереи исходных изображений, отображение результатов обнаружения дефектов дорожного покрытия, обновление сетки результатов при изменении размеров окна и управление отображением панели входных изображений.

\subsubsection{Класс "<GalleryManager">}

Методы класса "<GalleryManager"> представлены в таблице \ref{gallerymanager:table}.

\renewcommand{\arraystretch}{0.8}
\begin{xltabular}{\textwidth}{|p{4.5cm}|p{3.5cm}|>{\setlength{\baselineskip}{0.7\baselineskip}}X|}
	\caption{Описание методов класса "<GalleryManager">\label{gallerymanager:table}}\\
	\hline \centrow \setlength{\baselineskip}{0.7\baselineskip} Название метода & \centrow \setlength{\baselineskip}{0.7\baselineskip} Аргументы & \centrow Описание метода \\
	\hline \centrow 1 & \centrow 2 & \centrow 3\\ \hline
	\endfirsthead
	
	\caption*{Продолжение таблицы \ref{gallerymanager:table}}\\
	\hline \centrow 1 & \centrow 2 & \centrow 3\\ \hline
	\finishhead
	
	\_\_init\_\_ & window &
	Конструктор класса. Сохраняет ссылку на главное окно приложения для взаимодействия с элементами интерфейса. \\ \hline
	
	update\_result\_grid & Аргументов нет &
	Формирование и обновление сетки обработанных изображений. Выполняет размещение миниатюр результатов обработки и информации о количестве обнаруженных объектов. \\ \hline
	
	display\_results & Аргументов нет &
	Загрузка обработанных изображений, создание их уменьшенных копий и подготовка данных для отображения в галерее результатов. \\ \hline
	
	show\_thumbnails & Аргументов нет &
	Загрузка и отображение миниатюр исходных изображений в левой панели интерфейса. Обеспечивает отображение индикатора выполнения процесса загрузки. \\ \hline
	
	clear\_layout & layout &
	Очистка элементов пользовательского интерфейса, расположенных в контейнере \texttt{layout}. \\ \hline
	
	toggle\_input\_visibility & checked &
	Управление отображением панели исходных изображений в зависимости от значения параметра \texttt{checked}. \\ \hline
	
\end{xltabular}
\renewcommand{\arraystretch}{1.0}

\subsection{Модуль "<viewer.py">}

Модуль "<viewer.py"> содержит класс "<ImageViewer">, предназначенный для просмотра изображений в отдельном окне. Модуль обеспечивает отображение выбранного изображения, масштабирование изображения при изменении размеров окна и навигацию между изображениями.

\subsubsection{Класс "<ImageViewer">}

Методы класса "<ImageViewer"> представлены в таблице \ref{viewer:table}.

\renewcommand{\arraystretch}{0.8}
\begin{xltabular}{\textwidth}{|p{3.5cm}|p{3.5cm}|>{\setlength{\baselineskip}{0.7\baselineskip}}X|}
	\caption{Описание методов класса "<ImageViewer">\label{viewer:table}}\\
	\hline \centrow \setlength{\baselineskip}{0.7\baselineskip} Название метода & \centrow \setlength{\baselineskip}{0.7\baselineskip} Аргументы & \centrow Описание метода \\
	\hline \centrow 1 & \centrow 2 & \centrow 3\\ \hline
	\endfirsthead
	
	\caption*{Продолжение таблицы \ref{viewer:table}}\\
	\hline \centrow 1 & \centrow 2 & \centrow 3\\ \hline
	\finishhead
	
	\_\_init\_\_ & images, index &
	Конструктор класса. Выполняет создание окна просмотра изображений, настройку элементов интерфейса и отображение изображения с номером \texttt{index} из списка \texttt{images}. \\ \hline
	
	show\_image & Аргументов нет &
	Загрузка и отображение текущего изображения с автоматическим масштабированием под размеры окна просмотра. \\ \hline
	
	resizeEvent & event &
	Обработчик изменения размеров окна. Выполняет повторное масштабирование и отображение изображения после изменения размеров окна просмотра. \\ \hline
	
	prev\_image & Аргументов нет &
	Переход к предыдущему изображению в списке изображений. \\ \hline
	
	next\_image & Аргументов нет &
	Переход к следующему изображению в списке изображений. \\ \hline
	
\end{xltabular}
\renewcommand{\arraystretch}{1.0}

\subsection{Модуль "<theme\_manager.py">}

Модуль "<theme\_manager.py"> содержит класс "<ThemeManager">, предназначенный для управления визуальным оформлением пользовательского интерфейса. Модуль обеспечивает переключение между светлой и тёмной темами оформления и сохранение выбранных параметров интерфейса.

\subsubsection{Класс "<ThemeManager">}

Методы класса "<ThemeManager"> представлены в таблице \ref{theme:table}.

\renewcommand{\arraystretch}{0.8}
\begin{xltabular}{\textwidth}{|p{4.0cm}|p{3.5cm}|>{\setlength{\baselineskip}{0.7\baselineskip}}X|}
	\caption{Описание методов класса "<ThemeManager">\label{theme:table}}\\
	\hline \centrow \setlength{\baselineskip}{0.7\baselineskip} Название метода & \centrow \setlength{\baselineskip}{0.7\baselineskip} Аргументы & \centrow Описание метода \\
	\hline \centrow 1 & \centrow 2 & \centrow 3\\ \hline
	\endfirsthead
	
	\caption*{Продолжение таблицы \ref{theme:table}}\\
	\hline \centrow 1 & \centrow 2 & \centrow 3\\ \hline
	\finishhead
	
	\_\_init\_\_ & window &
	Конструктор класса. Сохраняет ссылку на главное окно приложения. \\ \hline
	
	set\_theme & theme &
	Установка темы оформления интерфейса в соответствии со значением параметра \texttt{theme}. Выполняет сохранение выбранной темы в настройках приложения. \\ \hline
	
	apply\_dark\_theme & Аргументов нет &
	Применение тёмной темы оформления пользовательского интерфейса. \\ \hline
	
	apply\_light\_theme & Аргументов нет &
	Применение светлой темы оформления пользовательского интерфейса. \\ \hline
	
\end{xltabular}
\renewcommand{\arraystretch}{1.0}


\subsection{Пользовательский интерфейс программы}

Пользовательский интерфейс программно-информационной системы обнаружения дефектов дорожного покрытия реализован с использованием библиотеки PySide6 и построен по принципу графического оконного приложения.

После запуска программы пользователю отображается стартовый экран с краткой инструкцией по работе с системой и кнопками быстрого доступа к основным функциям (рисунок \ref{startscr:image}). Для начала работы необходимо выбрать изображения для анализа с помощью команды открытия файлов.

\begin{figure}[H]
	\center{\includegraphics[width=1\linewidth]{startscr}}
	\caption{Стартовый экран}
	\label{startscr:image}
\end{figure}

В верхней части окна расположена строка меню, содержащая разделы «Файл», «Обработка» и «Вид». Раздел «Файл» предназначен для загрузки изображений, сохранения результатов обработки и открытия ранее сохранённых результатов. Раздел «Обработка» позволяет запускать процесс обнаружения дефектов дорожного покрытия, выбирать порог уверенности нейросетевой модели и управлять каталогом временных данных. Раздел «Вид» предоставляет средства управления отображением исходных изображений, сортировкой результатов и выбором темы оформления интерфейса.

В левой части главного окна отображается список загруженных изображений в виде миниатюр. Выбор миниатюры позволяет открыть соответствующее изображение в отдельном окне просмотра. Для уменьшения занимаемого пространства панель исходных изображений может быть скрыта пользователем.

Правая часть окна предназначена для отображения результатов обработки. После завершения обнаружения дефектов дорожного покрытия обработанные изображения отображаются в виде адаптивной сетки. Для каждого изображения дополнительно выводится количество обнаруженных дефектов. При изменении размеров окна расположение элементов автоматически перестраивается.

В нижней части окна располагается строка состояния, содержащая информацию о количестве загруженных и обработанных изображений, общем числе обнаруженных дефектов, среднем количестве дефектов на изображение, времени выполнения обработки и текущем состоянии программы. Также отображается индикатор выполнения операций загрузки и обработки изображений.

Для удобства работы предусмотрено отдельное окно просмотра изображений, позволяющее просматривать исходные и обработанные изображения в увеличенном масштабе. Окно просмотра поддерживает переход к предыдущему и следующему изображениям без необходимости возврата к главному окну программы.

В программе реализованы светлая и тёмная темы оформления интерфейса. Выбранная пользователем тема автоматически сохраняется и восстанавливается при последующих запусках приложения.

\subsection{Тестирование разработанной программной системы}

Для проверки работоспособности программно-информационной системы обнаружения дефектов дорожного покрытия было выполнено комплексное тестирование, включающее модульное, интеграционное и системное тестирование. Проверка осуществлялась на наборах изображений дорожного покрытия с различным разрешением и содержанием дефектов.

\subsubsection{Модульное тестирование компонентов системы}

Модульное тестирование выполнялось с целью проверки корректности работы отдельных компонентов программной системы. Для автоматизации проверки использовался фреймворк \texttt{pytest}.

Для класса \texttt{Detector} проверялась корректность изменения порога уверенности обнаружения объектов, а также корректность выполнения метода \texttt{process\_images()}, включая формирование структуры возвращаемых данных после обработки изображений.

Для класса \texttt{FileManager} проверялась работа механизма формирования текстового отчёта, а также корректность загрузки и анализа ранее сохранённого файла отчёта.

Результаты модульного тестирования представлены на рисунке \ref{Testmod:image}.

\begin{figure}[H]
	\begin{lstlisting}[language={}]
		tests/test_detector.py::test_set_conf PASSED
		tests/test_detector.py::test_process_images_basic PASSED
		tests/test_file_manager.py::test_create_report PASSED
		tests/test_report_parser.py::test_load_report PASSED
	\end{lstlisting}  
	\caption{Результат модульного тестирования}
	\label{Testmod:image}
\end{figure}

\subsubsection{Системное тестирование}

Системное тестирование выполнялось для проверки работоспособности программно-информационной системы обнаружения дефектов дорожного покрытия в целом. В ходе тестирования проверялись основные пользовательские сценарии работы.

Сценарий 1: загрузка изображений. Проверялась возможность загрузки одного или нескольких изображений с использованием диалога выбора файлов. После выбора файлов система должна отображать миниатюры изображений в левой части главного окна и обновлять информацию о количестве загруженных изображений.

В результате тестирования установлено, что изображения корректно загружаются в программу, отображаются в виде миниатюр.

Результат загрузки изображений представлен на рисунке \ref{load_images:image}.

\begin{figure}[H]
	\center{\includegraphics[width=1\linewidth]{load_images}}
	\caption{Загрузка изображений в программу}
	\label{load_images:image}
\end{figure}

Сценарий 2: настройка параметров обнаружения. Проверялась возможность изменения порога уверенности нейросетевой модели через меню программы. Для проверки последовательно выбирались все доступные значения порога обнаружения.

В ходе тестирования было установлено, что выбранное значение корректно передаётся модулю обнаружения объектов и используется при последующей обработке изображений.

Настройка параметров обнаружения представлена на рисунке \ref{conf:image}.

\begin{figure}[H]
	\center{\includegraphics[width=1\linewidth]{conf}}
	\caption{Настройка порога обнаружения}
	\label{conf:image}
\end{figure}

Сценарий 3: анализ изображений. Проверялась работа алгоритма обнаружения дефектов дорожного покрытия на наборе тестовых изображений. В процессе тестирования выполнялась загрузка изображений различных форматов, их автоматическое масштабирование при превышении допустимого размера и выполнение анализа с использованием обученной модели YOLOv8n.

Также внимание уделялось корректности формирования ограничивающих рамок обнаруженных объектов, отображению подписи класса и значений уверенности модели. Дополнительно проверялись подсчёт общего количества обнаруженных дефектов, подсчёт количества дефектов на каждом изображении и обновление другой статистической информации в строке состояния.

В ходе тестирования подтверждена корректная работа алгоритма обнаружения дефектов дорожного покрытия, а также правильность формирования результатов обработки.

В результате тестирования было подтверждено корректное выполнение обнаружения дефектов дорожного покрытия, подсчёт найденных объектов и сохранение результатов обработки во временном каталоге.

Результат анализа изображений представлен на рисунке \ref{detection:image}.

\begin{figure}[H]
	\center{\includegraphics[width=1\linewidth]{detection}}
	\caption{Обнаружение дефектов дорожного покрытия}
	\label{detection:image}
\end{figure}

Сценарий 4: просмотр результатов обнаружения. Проверялась возможность просмотра результатов обработки в главном окне программы и в отдельном окне просмотра изображений, тестировалась навигация между изображениями с использованием кнопок перехода.

В ходе тестирования установлено, что результаты корректно отображаются в виде галереи, а просмотр изображений осуществляется без ошибок.

Просмотр результатов обнаружения представлен на рисунке \ref{gallery:image}.

\begin{figure}[H]
	\center{\includegraphics[width=1\linewidth]{gallery}}
	\caption{Просмотр результатов обработки}
	\label{gallery:image}
\end{figure}

Сценарий 5: сортировка результатов. Проверялась работа механизмов сортировки обработанных изображений по имени файла и по количеству обнаруженных объектов.

Результат сортировки представлен на рисунке \ref{sort:image}.

\begin{figure}[H]
	\center{\includegraphics[width=1\linewidth]{sort}}
	\caption{Сортировка результатов обработки}
	\label{sort:image}
\end{figure}

Сценарий 6: Сохранение результатов и формирование отчёта. Проверялась возможность сохранения обработанных изображений и автоматического формирования текстового отчёта о результатах анализа. В процессе тестирования контролировалась корректность копирования файлов в выбранный пользователем каталог, а также работа механизма обработки конфликтов имён существующих файлов.

Дополнительно проверялось формирование файла \texttt{report.txt}, содержащего количество обработанных изображений, число обнаруженных дефектов дорожного покрытия, среднее количество объектов на изображение, время обработки и статистику по каждому обработанному изображению.

В ходе тестирования было установлено, что программа корректно сохраняет изображения в выбранный каталог, формирует файл \texttt{report.txt} и записывает в него статистическую информацию о выполненной обработке.

Результат сохранения данных представлен на рисунке \ref{saved:image}.

\begin{figure}[H]
	\center{\includegraphics[width=1\linewidth]{saved}}
	\caption{Сохранение результатов обработки}
	\label{saved:image}
\end{figure}

Сценарий 7: загрузка ранее сохранённых результатов. Проверялась возможность открытия для просмотра ранее сохранённых результатов обработки. Контролировалось восстановление статистической информации о выполненной обработке, загрузка обработанных изображений и формирование галереи результатов. Для этого использовался каталог, содержащий обработанные изображения и файл отчёта. 

В результате тестирования установлено, что программа корректно восстанавливает статистические данные, загружает список изображений и отображает результаты обработки.

Загрузка ранее сохранённых результатов представлена на рисунке \ref{open_results:image}.

\begin{figure}[H]
	\center{\includegraphics[width=1\linewidth]{open_results}}
	\caption{Загрузка ранее сохранённых результатов}
	\label{open_results:image}
\end{figure}

Сценарий 8: очистка временных данных. Проверялась работа механизма очистки временного каталога, используемого для хранения промежуточных результатов обработки.

В ходе тестирования было подтверждено корректное удаление временных файлов и последующее создание пустого каталога для дальнейшей работы программы.

Очистка временных данных представлена на рисунках \ref{clear_bef:image} и \ref{clear_aft:image}.

\begin{figure}[H]
	\center{\includegraphics[width=1\linewidth]{clear_bef}}
	\caption{До очистки временных данных}
	\label{clear_bef:image}
\end{figure}

\begin{figure}[H]
	\center{\includegraphics[width=1\linewidth]{clear_aft}}
	\caption{После очистки временных данных}
	\label{clear_aft:image}
\end{figure}

Сценарий 9: изменение директории сохранения временных файлов. Проверялась возможность выбора пользователем другой директории для сохранения временных файлов результатов обработки изображений.

В результате тестирования установлено, что выбранная пользователем директория корректно сохраняется в настройках приложения и используется при следующей обработке для временного хранения файлов.

Результаты изменения директории сохранения представлены на рисунке \ref{change_dir:image}.

\begin{figure}[H]
	\center{\includegraphics[width=1\linewidth]{change_dir}}
	\caption{Изменение директории сохранения файлов}
	\label{change_dir:image}
\end{figure}

Сценарий 10: переключение темы интерфейса. Проверялась возможность изменения оформления пользовательского интерфейса между светлой и тёмной темами.

Результаты тестирования показали, что тема оформления корректно изменяется без перезапуска программы, а выбранные настройки сохраняются и автоматически восстанавливаются при следующем запуске приложения.

Переключение темы интерфейса представлено на рисунке \ref{theme:image}.

\begin{figure}[H]
	\center{\includegraphics[width=1\linewidth]{theme}}
	\caption{Переключение темы интерфейса}
	\label{theme:image}
\end{figure}

Тестирование пользовательского интерфейса. Проверка интерфейсных компонентов выполнялась вручную.

Были протестированы:
\begin{itemize}
	\item отображение миниатюр входных изображений;
	\item отображение миниатюр результатов обработки;
	\item работа прогресс-бара;
	\item навигация между изображениями;
	\item адаптивное изменение сетки результатов.
\end{itemize}

Критических нарушений в работе интерфейса не было выявлено.











\begin{comment}
\subsection{Классы, используемые при разработке сайта}

Можно выделить следующий список классов и их методов, использованных при разработке web-приложения (таблица \ref{class:table}). Пример таблицы с уменьшенным межстрочным интервалом.

\renewcommand{\arraystretch}{0.8} % уменьшение расстояний до сетки таблицы
\begin{xltabular}{\textwidth}{|X|p{2.5cm}|>{\setlength{\baselineskip}{0.7\baselineskip}}p{4.85cm}|>{\setlength{\baselineskip}{0.7\baselineskip}}p{4.85cm}|}
\caption{Описание классов Bitrix, используемых в приложении\label{class:table}}\\
\hline \centrow \setlength{\baselineskip}{0.7\baselineskip} Название класса & \centrow \setlength{\baselineskip}{0.7\baselineskip} Модуль, к которому относится класс & \centrow Описание класса & \centrow Методы \\
\hline \centrow 1 & \centrow 2 & \centrow 3 & \centrow 4\\ \hline
\endfirsthead
\caption*{Продолжение таблицы \ref{class:table}}\\
\hline \centrow 1 & \centrow 2 & \centrow 3 & \centrow 4\\ \hline
\finishhead
CMain & Главный модуль & CMain – главный класс страницы web-приложения. После одного из этапов по загрузке страницы в сценарии становится доступным инициализированный системой объект данного класса с именем \$APPLICATION & void ShowTitle(string property\_code = «title», bool strip\_tags = true)
Выводит заголовок страницы
void SetTitle(string title)
Устанавливает заголовок страницы

void ShowCSS(bool external = true, bool XhtmlStyle = true)
Выводит таблицу стилей CSS страницы\\
\hline CFile & Главный модуль & CFile – Класс для работы с файлами и изображениями & array GetFileArray (int file\_id)
Метод возвращает массив, содержащий описание файла (путь к файлу, имя файла, размер) с идентификатором file\_id
\end{xltabular}
\renewcommand{\arraystretch}{1.0} % восстановление сетки

\subsection{Модульное тестирование разработанного web-сайта}

Модульный тест для класса User из модели данных представлен на рисунке \ref{unitUser:image}.

\begin{figure}[ht]
\begin{lstlisting}[language=Python]
from django.test import TestCase
from .models import *
User = get_user_model()


class ShpoTestCases(TestCase):

    def setUp(self) -> None:
        self.user = User.objects.create(username='testtestovich', password='testtestovich', first_name='Sad', last_name='')

    def test_2(self):

        self.assertEqual(self.user.first_name, 'Sad')
        self.assertEqual(self.user.last_name, 'Cat')
        print((self.user))
        print((self.user.first_name))
        print((self.user.last_name))
\end{lstlisting}  
\caption{Модульный тест класса User}
\label{unitUser:image}
\end{figure}

\subsection{Системное тестирование разработанного web-сайта}

На рисунке \ref{main:image} представлена главная страница сайта «Русатом – Аддитивные технологии».
\newpage % при необходимости можно переносить рисунок на новую страницу
\begin{figure}[H] % H - рисунок обязательно здесь, или переносится, оставляя пустоту
\center{\includegraphics[width=1\linewidth]{main1}}
\center{\includegraphics[width=1\linewidth]{main2}}
\center{\includegraphics[width=1\linewidth]{main3}}
\caption{Главная страница сайта «Русатом – Аддитивные технологии»}
\label{main:image}
\end{figure}

На рисунке \ref{menu:image} представлен динамический вывод заголовков, включающий в себя искомые фразы при поиске фраз.

\begin{figure}[ht]
\center{\includegraphics[width=1\linewidth]{menu}}
\caption{Динамический вывод заголовков}
\label{menu:image}
\end{figure}

На рисунке \ref{enter:image} представлен ввод данных для публикации новости.

\begin{figure}[ht]
\center{\includegraphics[width=1\linewidth]{enter}}
\caption{Ввод данных для публикации очень-очень длинной, интересной и полезной новости}
\label{enter:image}
\end{figure}
\end{comment}
